% Research draft for R75; deliberately outside the sealed R74 source.
% Depends on the retained permutation-core normal form and positivity argument.
\section{Orbit multiplicity, positive synthesis, and the convexity boundary}
\label{sec:r75-orbit-multiplicity}

The quotient reduction of Section~\ref{sec:r74-quotient} has a precise
boundary. A single transitive orbit admits convexification without enlarging
the orbit exactly when its stabilizer is normal. For general finite groups,
irreducible multiplicities measure the obstruction and determine the exact
number of prescribed orbit copies at sufficiently small positive contrast.
All representations below are finite dimensional. We use classical complete
reducibility, matrix-coefficient orthogonality, and Frobenius reciprocity
\cite{Serre77}; their representation-theoretic content is not claimed as new.

Let $G$ be a finite group, not necessarily the faithful observable group of
the particular table under discussion. For a nonempty finite right $G$-set
$X$, let $\mathcal R_X(C)$ denote the probability tables
\[
 b(g)=\lambda P_gD,\qquad
 \lambda\in\Delta(X),\quad D(x,\cdot)\in\Delta(C),
\]
where $P_gP_h=P_{gh}$. Thus the same preparation and decoder are used for
every group element. Taking an observable quotient does not change the
table; keeping $G$ explicit lets us compare different allowed actions.
Write $\widehat G$ for the complex irreducible representations, choose a
unitary representative $\rho$ of degree $d_\rho$, and define
\begin{equation}\label{eq:r75-fourier-block}
 \widehat b_c(\rho)=\frac1{|G|}\sum_{g\in G}b(g,c)\rho(g^{-1}),
 \qquad
 A_\rho(b)=\begin{pmatrix}
 \widehat b_{c_1}(\rho)\\ \vdots\\ \widehat b_{c_{|C|}}(\rho)
 \end{pmatrix},
 \qquad r_\rho(b)=\operatorname{rank}_{\mathbb C}A_\rho(b).
\end{equation}
Changing unitary bases does not change these ranks.

\begin{lemma}[Fourier multiplicity and real descent]
\label{lem:r75-fourier-multiplicity}
Let $V_b$ be the real span of the translates $b^h(g)=b(hg)$. The
multiplicity of the right irreducible $\rho$ in $V_b\otimes_{\mathbb R}
\mathbb C$ is $r_\rho(b)$. In particular $r_{\boldsymbol 1}(b)=1$ and
$r_\rho(b)\le d_\rho$. If $U$ and $V$ are real $G$-modules, then $U$
embeds equivariantly in $V$ if and only if every complex irreducible
multiplicity in $U\otimes\mathbb C$ is at most the corresponding one in
$V\otimes\mathbb C$.
\end{lemma}
\begin{proof}
The Fourier transform sends $b^h$ to the blocks
$A_\rho(b)\rho(h)$. Matrix-coefficient orthogonality identifies
$\mathbb C[G]$ with $\bigoplus_{\rho\in\widehat G}M_{d_\rho}(\mathbb C)$:
the group algebra can independently prescribe an arbitrary matrix on each
irreducible block. Consequently the $\rho$-part of the cyclic module
generated by $b$ identifies with
\[
 \{A_\rho(b)M:M\in M_{d_\rho}(\mathbb C)\}.
\]
If $A_\rho$ has rank $r$, a choice of a basis of its column space
identifies this right module with $r$ independent rows of length
$d_\rho$. It therefore consists of $r$ copies of $\rho$. For the trivial
representation the stacked block is the nonzero probability vector
$\pi=|G|^{-1}\sum_gb(g)$, so its rank is one.

Necessity of the second assertion follows by complexification. For
sufficiency, complete reducibility over $\mathbb C$ and the multiplicity
inequalities give a complex equivariant injection
$F:U\otimes\mathbb C\to V\otimes\mathbb C$. In real bases write
$F=A+iB$. Both $A$ and $B$ are real intertwiners. A maximal nonzero minor
of $F$ is the value at $z=i$ of a real polynomial minor of $A+zB$.
That polynomial is not identically zero, so some real $z$ makes this
minor nonzero. The real map $A+zB$ is the required injection. One may
choose $z$ among $0,1,\ldots,\dim U$, since the nonzero polynomial has
degree at most $\dim U$. This proves the assertion without excluding
complex or quaternionic representation types.
\end{proof}

For $H\le G$ put
\[
 m_\rho(H)=\dim_{\mathbb C}\rho^H
 =\operatorname{rank}\left(\frac1{|H|}\sum_{h\in H}\rho(h)\right).
\]
By Frobenius reciprocity these are the multiplicities of
$\mathbb C[H\backslash G]$. Let $q_H(b)$ be the least positive integer
$q$ for which $b$ belongs to $\mathcal R_{q(H\backslash G)}(C)$, where
$q(H\backslash G)$ is the disjoint union of $q$ copies; set $q_H(b)=\infty$
if no such factorization exists. This is an action-restricted resource.
It is not $\kappa_G(b)/[G:H]$ in general, since an unrestricted optimal
core may use other orbit types.

\begin{theorem}[Exact orbit multiplicity at small contrast]
\label{thm:r75-orbit-multiplicity}
Let $b:G\to\Delta(C)$, let
$\pi=|G|^{-1}\sum_gb(g)$, and set $b_t=\pi+t(b-\pi)$. For $H\le G$ define
\begin{equation}\label{eq:r75-copy-number}
 Q_H(b)=\max_{\rho\in\widehat G}
 \left\lceil\frac{r_\rho(b)}{m_\rho(H)}\right\rceil,
\end{equation}
with a positive numerator and zero denominator interpreted as $\infty$
and $0/0$ as zero. For every $t>0$,
$q_H(b_t)\ge Q_H(b)$. If $Q_H(b)<\infty$, an explicitly positive $t_0$
exists such that
\[
 q_H(b_t)=Q_H(b)\qquad(0<t\le t_0).
\]
If $Q_H(b)=\infty$, no positive contrast is realizable using any finite
number of copies of this orbit. At zero contrast one copy suffices.

More generally, choose one representative of each subgroup conjugacy
class. The unrestricted permutation-module bound has the exact finite
integer formulation
\begin{equation}\label{eq:r75-multiplicity-cover}
 \nu_G(V_b)=
 \min\left\{
  \sum_H q_H[G:H]:\ q_H\in\mathbb Z_{\ge0},\quad
  \sum_Hq_Hm_\rho(H)\ge r_\rho(b)\quad(\rho\in\widehat G)
 \right\}.
\end{equation}
Consequently this integer minimum equals $\kappa_G(b_t)$ for every
sufficiently small positive contrast, with the pure-start charge given
by the retained normal-form theorem when $G$ is the observable group.
\end{theorem}
\begin{proof}
A factorization through a finite $G$-set $Y$ makes $V_{b_t}$ a quotient
of a submodule of $\mathbb R^Y$. Averaging an inner product over the
finite group makes every invariant subspace have an invariant orthogonal
complement; hence this subquotient embeds into $\mathbb R^Y$. The
contrast change is the identity on the invariant line generated by
$\pi$ and multiplication by $t$ on the zero-average part, so for $t>0$
its module is isomorphic to $V_b$. On $q(H\backslash G)$ the available
complex multiplicities are $qm_\rho(H)$, proving the lower bound and
the assertion when an active type is absent.

For the matching construction, let $Y$ be any nonempty finite $G$-set
whose complex multiplicities dominate $r_\rho(b)$. The lemma gives a
real equivariant injection of $V_b$ into $\mathbb R^Y$. Here are the
positivity details, to make clear that no complex factorization is
being substituted for a stochastic one. Delete the output columns
with $\pi_c=0$; nonnegativity makes them identically zero, and they can
be restored afterwards. Decompose
\[
 V_b=\mathbb R\overline b\oplus W,\qquad
 \overline b(g)=\pi,
\]
where $W$ is the kernel of group averaging. Every element of $W$ has
zero total output mass. Put $k=|Y|$, let $p$ be its uniform probability
row, and choose an equivariant injection $j:W\to\mathbb R^Y$.
The image of $j$ is orthogonal to all invariant vectors. Let
$F:\mathbb R^Y\to W$ be orthogonal projection onto $j(W)$ followed by
$j^{-1}$. Then $Fj=I_W$ and $F(p)=0$. Set
\[
 a=j(b-\overline b),\qquad T(y,c)=F(e_y)(e,c).
\]
Equivariance gives $aP_gT=b(g)-\pi$, while $pT=0$ and each row of $T$
sums to zero. If $b$ is nonconstant, choose
\[
 \epsilon=\frac1{2k\|a\|_\infty},\qquad
 B=\max_{y,c}\frac{|T(y,c)|}{\pi_c},\qquad
 t_0=\min\{1,\epsilon/B\}.
\]
The probability row $\lambda=p+\epsilon a$ and the row-stochastic
decoder $D_t(y,c)=\pi_c+(t/\epsilon)T(y,c)$ then satisfy
$\lambda P_gD_t=b_t(g)$ for $0\le t\le t_0$. For a constant table the
constant decoder works at every contrast. This is precisely the
retained positive construction, applied to the specified action
rather than only to a globally minimum one. Taking
$Y=Q_H(b)(H\backslash G)$ proves equality in the copy formula.

Every finite $G$-set is a disjoint union of coset orbits, and each such
orbit contributes the multiplicities $m_\rho(H)$. The real descent
lemma therefore makes the displayed integer constraints equivalent to
real module containment. Conversely every feasible integer vector
defines that disjoint-union action. This proves
\eqref{eq:r75-multiplicity-cover}; the retained small-contrast theorem
or the preceding construction gives its positive realization. The
regular orbit makes the minimum finite, since
$m_\rho(\{e\})=d_\rho\ge r_\rho(b)$.
\end{proof}

\begin{theorem}[Normality is the exact single-orbit convexity boundary]
\label{thm:r75-normality-convexity}
For $H\le G$ and the transitive right action $X=H\backslash G$, the
following conditions are equivalent:
\begin{enumerate}
\item $H$ is normal in $G$;
\item $\mathcal R_X(\{0,1\})$ is convex;
\item $\mathcal R_X(C)$ is convex for every finite outcome set $C$.
\end{enumerate}
If $H$ is normal, this response set is exactly the polytope of
probability tables constant on the left cosets of $H$. If $H$ is not
normal, nonconvexity occurs in every neighborhood of the interior
constant binary table $(1/d,1-1/d)$, where $d=[G:H]$.

An explicit witness is supplied by the return tables
\[
 f_x(g)=\mathbf1_{\{xg=x\}},\qquad x\in X,
 \qquad
 F_X(g)=\frac1d\sum_{x\in X}f_x(g)=\frac{\chi_X(g)}d.
\]
Every $(f_x,1-f_x)$ lies in $\mathcal R_X(\{0,1\})$, whereas
$(F_X,1-F_X)$ does not when $H$ is not normal. The same conclusion
holds after replacing every scalar $f$ by $d^{-1}+t(f-d^{-1})$, for
every $0<t\le1$.
\end{theorem}
\begin{proof}
If $H$ is normal, for
$b(g,c)=\sum_{Hx}\lambda(Hx)D(Hxg,c)$ one has
$b(hg,c)=b(g,c)$ for every $h\in H$: normality gives $Hxh=Hx$.
Conversely every probability table with this left invariance is
realized by the point preparation $H$ and decoder
$D(Hg,c)=b(g,c)$. This proves the stated polytope identity and
(i)$\Rightarrow$(iii)$\Rightarrow$(ii).

Suppose $H$ is not normal. The group-algebra element
$e_H=|H|^{-1}\sum_{h\in H}h$ is an orthogonal projection in every
unitary irreducible representation, of rank $m_\rho(H)$. If every one
of these ranks were either zero or $d_\rho$, each matrix $\rho(e_H)$
would be zero or the identity. The faithful decomposition of the
complex group algebra into its irreducible matrix blocks would then
make $e_H$ central. Its coefficients are constant under conjugation
exactly when $H$ is normal. Hence some irreducible $\rho$ has
\[
 0<m_\rho(H)<d_\rho.
\]
This irreducible is nontrivial.

The return table at $x$ is realized by initial point $x$ and the
indicator decoder of $x$. Since
$\chi_X=\sum_\sigma m_\sigma(H)\chi_\sigma$, orthogonality gives
\[
 \widehat F_X(\rho)=\frac{m_\rho(H)}{d\,d_\rho}I_{d_\rho}.
\]
Its rank is $d_\rho>m_\rho(H)$. Lemma~\ref{lem:r75-fourier-multiplicity}
and the module lower bound therefore exclude a single-orbit
factorization. This average of members of the response set proves
nonconvexity. For positive contrast, the nontrivial Fourier block is
multiplied by $t$ and keeps its rank. Each contrasted return table is
still realized on its original orbit by replacing its decoder with
$d^{-1}+t(D-d^{-1})$. Letting $t$ decrease to zero proves the
neighborhood assertion. Thus (ii) implies (i).
\end{proof}

For this witness the exact small-contrast number of copies is
\begin{equation}\label{eq:r75-character-copies}
 Q_H(F_X,1-F_X)
 =\max_{\rho:m_\rho(H)>0}
   \left\lceil\frac{d_\rho}{m_\rho(H)}\right\rceil.
\end{equation}
Indeed every active Fourier block of the character is full rank,
including the trivial block. Thus nonnormality entails a strict copy
cost at arbitrarily small positive contrast, with a matching positive
realization. This obstruction is not caused by a decoder touching the
boundary of its simplex.

\begin{lemma}[Full-contrast return-law rigidity]
\label{lem:r75-return-rigidity}
Let $X$ be a transitive $d$-point $G$-set. If a scalar response is a
convex combination of responses on copies of $X$, has value one at
$e$, and has group average $1/d$, every positive-weight component
uses a point preparation and its point indicator decoder. In
particular every component is a return law $f_x$.
\end{lemma}
\begin{proof}
For a component write $\lambda\in\Delta(X)$ and $D\in[0,1]^X$.
The total value one at $e$ forces $\lambda D=1$ for every component
of positive weight, hence $D(x)=1$ on the nonempty support of
$\lambda$. Transitivity gives its group average
$d^{-1}\sum_xD(x)\ge1/d$. Equality of the weighted mean with $1/d$
forces equality for each component. Consequently $\sum_xD(x)=1$,
so the support of $\lambda$ consists of exactly one point and $D$
is its indicator.
\end{proof}

These statements separate three matters: the classical multiplicities
of a permutation representation, the positivity of its preparation
and decoder, and the cost of an externally prescribed pure start.
The integer formula is exact at small positive contrast; it does not
replace the general full-contrast bilinear synthesis problem by a
linear program. The convexity theorem pinpoints why the retained
normal-stabilizer absorption cannot simply be extended to arbitrary
subgroups.
